\documentclass[tikz,border=7pt]{standalone}
% flashcards-title: Two rays with a common endpoint
% flashcards-desc: Rays from B pass through A and C, and a small curved mark shows the angle between them without naming it.
\input{figures/tikz-style.tex}
\begin{document}
\begin{tikzpicture}[x=1cm,y=1cm]
  \coordinate (B) at (-1.8,-1.2);
  \coordinate (A) at (2.6,-1.2);
  \coordinate (C) at (1.45,1.65);
  \draw[geom ray] (B)--(A);
  \draw[geom ray] (B)--(C);
  \node[geom point,label={[geom small]below left:$B$}] at (B) {};
  \node[geom point,label={[geom small]below:$A$}] at (1.35,-1.2) {};
  \node[geom point,label={[geom small]above left:$C$}] at (0.55,0.9) {};
  \draw[geom construction] ($(B)+(0:0.85)$) arc[start angle=0,end angle=41,radius=0.85];
\end{tikzpicture}
\end{document}
